\section{最短闭环：从模型文件到 logits}

tiny MLP 路径不是为了证明 LCQI 已经是完整大模型引擎，而是为了建立最短的、可手算的、可测试的推理闭环。读者要把模型加载、输入检查、int8 linear、ReLU、第二层 logits、argmax 和 CLI 输出连起来看。

\section{模型加载与基础推理}

\filepath{model.cpp} 负责解析 tiny 文本模型、权重、scale、bias 和 shape。坏格式、坏维度和坏数值都应该在这里被拒绝，不能拖到 kernel 内部才崩溃。

\lstinputlisting[language=C++,caption={src/model.cpp}]{../../../cpu-volume-3/labs/linux_cpu_inference/src/model.cpp}

\filepath{inference.cpp} 执行 int8 linear、ReLU、第二层 logits 和 predicted class。它是系统中最短的“模型资产到输出”路径。

\lstinputlisting[language=C++,caption={src/inference.cpp}]{../../../cpu-volume-3/labs/linux_cpu_inference/src/inference.cpp}

\filepath{main.cpp} 把模型路径和输入接到 \code{run\_inference}，并打印结果。CLI 应保持薄层，不复制模型解析和推理逻辑。

\lstinputlisting[language=C++,caption={src/main.cpp}]{../../../cpu-volume-3/labs/linux_cpu_inference/src/main.cpp}

\section{Int8 Kernel 基础实现}

\filepath{int8\_kernels.cpp} 包含 scalar、packed、workspace、deterministic fixture 和误差比较。它是 AVX2 路径的 reference 对照，也是 benchmark 的正确性底座。

\lstinputlisting[language=C++,caption={src/int8\_kernels.cpp}]{../../../cpu-volume-3/labs/linux_cpu_inference/src/int8_kernels.cpp}

\filepath{int8\_kernels\_avx2.cpp} 是平台优化路径。读者要先看可用性判断，再看向量化实现；如果机器不支持 AVX2/FMA，报告必须写 unavailable。

\lstinputlisting[language=C++,caption={src/int8\_kernels\_avx2.cpp}]{../../../cpu-volume-3/labs/linux_cpu_inference/src/int8_kernels_avx2.cpp}
